\section*{Limitations}
Despite its strong performance, StructMem has several limitations. The quality of dual-perspective extraction is highly dependent on instruction prompts, where suboptimal design may result in incomplete or inaccurate relational information capture. Future research could investigate automated prompt optimization to improve robustness across various dialogue contexts. Additionally, the framework primarily addresses memory expansion and synthesis but currently lacks an explicit mechanism for conflict resolution and memory updating. As user facts or preferences may evolve over long horizons, the absence of a revision process could lead to inconsistencies between historical summaries and new information. Future iterations should incorporate memory decay or updating strategies to ensure the hierarchical organization accurately reflects the most current state of the interaction.


\section*{Acknowledgements}
We would like to express sincere gratitude to the  reviewers for their thoughtful and constructive feedback. This work was supported by the National Natural Science Foundation of China (No. 62576307), Yongjiang Talent Introduction Programme (2021A-156-G), and Information Technology Center and State Key Lab of CAD\&CG, Zhejiang University. This work was supported by Ant Group and Zhejiang University - Ant Group Joint Laboratory of Knowledge Graph.